% Done
\section{Analýza požadavků}

% Review
V této sekci sepíšu funkční a nefunkční požadavky vznikající aplikace. Při sepisování požadavků budu vycházet zejména z analýzy stávajícího řešení a analýzy existujících aplikací. Dále budu také vycházet z výsledků uživatelských testů, které byly provedeny v rámci mé bakalářské práce \cite{Jenicek2024BachelorThesis} zabývající se stávající aplikací, a také z návrhu úprav do budoucna, které byly popsány ve zmíněné práci.

% Review
Po sepsání jednotlivých požadavků se zaměřím na určení jejich priorit pomocí metody MoSCoW \cite{Moscow}. Tato metoda umožňuje jednoduše rozdělit požadavky do čtyř kategorií dle jejich priority. Požadavky v kategorii \textit{must have} musí být při implementaci splněny, požadavky kategorie \textit{should have} by měly být splněny, požadavky kategorie \textit{could have} mohou být splněny, a požadavky kategorie \textit{won't have} nebudou v rámci této práce implementovány.

% Done
Na základě předchozí analýzy stávajícího řešení a analýzy konkurenčních aplikací jsem identifikoval a sepsal následující funkční a nefunkční požadavky.

% Done
\subsection{Funkční požadavky}

% Review
Při analýze jsem identifikoval a sepsal následující funkční požadavky, kterým se budu věnovat v rámci této práce.

% Done
\subsubsection{FP-1 Úvodní stránka aplikace}\label{sec:fp-welcome-page}

% Done
\begin{tabular}{lp{0.7\linewidth}}
    \textbf{Priorita:} & Should have \\
\end{tabular}

% Done
Aplikace by měla obsahovat úvodní stránku, která bude uživateli zobrazena, když není přihlášen, nebo když poprvé otevře aplikaci. Na této stránce by měl mít uživatel možnost zvolit, jestli se chce přihlásit nebo registrovat. Stránka by měla obsahovat logo a název aplikace a krátký popis hlavní hodnoty, kterou aplikace uživateli přinese.

% Done
Tento požadavek vychází z analýzy konkurence, kde většina aplikací obsahuje kromě stánky pro přihlášení a registraci i úvodní stránku s grafikou či textem, který prezentuje hlavní přínosy aplikace pro uživatele.

% Done
\subsubsection{FP-2 Nový úvodní dotazník a doporučování kurzů}\label{sec:fp-onboarding}

% Done
\begin{tabular}{lp{0.7\linewidth}}
    \textbf{Priorita:} & Should have \\
\end{tabular}

% Done
Aplikace by měla obsahovat úvodní dotazník zaměřený na získání informací o uživateli. Tyto informace by se měly týkat zejména cílů uživatele a jeho zkušeností s programováním. Na základě dotazníku by měla aplikace nabídnout uživateli několik kurzů, které odpovídají jeho cílům a zkušenostem.

% Done
Tato funkce je inspirována konkurenční aplikací Codecademy, která využívá úvodní dotazník pro doporučování kurzů uživateli.

% Done
\subsubsection{FP-3 Prohlížení a doporučování kurzů}\label{sec:fp-course-browse}

% Done
\begin{tabular}{lp{0.7\linewidth}}
    \textbf{Priorita:} & Must have \\
\end{tabular}

% Done
Uživatel by měl mít v aplikaci možnost zobrazit přehled dostupných kurzů a jednoduše je prohledávat. U každého kurzu by měl mít možnost se přesunout na stránku daného kurzu.

% Done
Součástí prohlížení kurzů by měla být sekce s doporučenými kurzy. Tato doporučení by měla být založena na úvodním dotazníku, který uživatele vyplnil při registraci. Uživatel by měl mít možnost dotazník kdykoliv vyplnit znovu a přizpůsobit si tak svoje preference při zapisování kurzů.

% Done
\subsubsection{FP-4 Stránka kurzu}\label{sec:fp-course-page}

% Done
\begin{tabular}{lp{0.7\linewidth}}
    \textbf{Priorita:} & Must have \\
\end{tabular}

% Done
Každý kurz v aplikaci by měl mít vlastní stránku s jeho popisem a přehledem kapitol, které jsou v kurzu obsaženy. Uživatel by měl mít možnost se do kurzu přihlásit. Dále by měl mít možnost se přesunout na stránku pro správu zapsaných kurzů, kde se bude moci z kurzu odhlásit. Pokud je uživatel do kurzu zapsaný, měl by mít možnost jednoduše pokračovat v učení daného kurzu. Uživatel by měl mít také možnost zobrazit svůj certifikát o dokončení kurzu, pokud už jej získal.

% Done
Tento požadavek vychází z analýzy konkurence, kde většina aplikací obsahuje stránky kurzů s jejich popisem a s přehledem lekcí, které jsou v kurzu obsaženy.

% Done
\subsubsection{FP-5 Stránka kapitoly kurzu}\label{sec:fp-course-chapter-page}

% Done
\begin{tabular}{lp{0.7\linewidth}}
    \textbf{Priorita:} & Must have \\
\end{tabular}

% Done
Aplikace by měla umožnit uživateli si zobrazit stránku kapitoly kurzu. Na této stránce by měly být zobrazeny všechny lekce dané kapitoly. Lekce by měly být seskupeny dle jejich témat do sekcí. Lekce by měly být označeny jako dokončené nebo nedokončené podle toho, zda uživatel danou lekci dokončil. Aplikace by také měla vizuálně zvýraznit lekci, kterou by se uživatel měl učit dále.

% Done
Ve stávající aplikaci jsou lekce kurzu vyobrazeny jako body na čáře připomínající cestu, po které se uživatel při postupování kurzem pohybuje. V rámci tohoto zobrazení tak nejsou zobrazeny názvy lekcí. Byť je toto zobrazení vizuálně poutavé, a například konkurenční aplikace Mimo toto zobrazení využívá, může být velice nepraktické zejména z důvodu, že je pro uživatele značně obtížné najít konkrétní lekci, pokud se ji chce naučit nebo si ji zopakovat.

% Done
Z hlediska použitelnosti jsem tak vyhodnotil, že by v rámci tohoto požadavku měl být navržen nový vzhled seznamu lekcí, který bude možné použít na stránce kapitoly kurzu. Tento vzhled by měl v základu fungovat jako seznam lekcí rozdělených do sekcí. Vizuálně by měl ale také částečně připomínat původní vizuální vyobrazení ve formě připomínající cestu s lekcemi. Cílem je tak vytvořit vizuálně poutavý vzhled znázorňující chronologickou posloupnost lekcí, který ovšem bude praktičtější a bude umožňovat uživateli se lépe v lekcích orientovat.

% Done
\subsubsection{FP-6 Certifikáty o dokončení kurzu}\label{sec:fp-course-certificate}

% Done
\begin{tabular}{lp{0.7\linewidth}}
    \textbf{Priorita:} & Should have \\
\end{tabular}

% Done
Po dokončení kurzu by měl uživateli s premium účtem být vygenerován certifikát o dokončení daného kurzu. Certifikát by měl být zobrazen na veřejně dostupné stránce a měl by obsahovat základní informace o dokončeném kurzu a uživateli, který kurz dokončil. Uživatel by měl mít možnost odkaz na certifikát jednoduše sdílet s ostatními uživateli. Uživatelům, kteří nemají premium účet, by měla být zobrazena informace, že si budou moci certifikát vygenerovat po zakoupení premium předplatného.

% Done
\subsubsection{FP-7 Správa zapsaných kurzů}\label{sec:fp-course-management}

% Done
\begin{tabular}{lp{0.7\linewidth}}
    \textbf{Priorita:} & Must have \\
\end{tabular}

% Done
Uživatel by měl mít možnost si zobrazit přehled zapsaných kurzů. U každého kurzu by měl uživatel mít možnost se z něj odhlásit.

% Done
\subsubsection{FP-8 Cesty učení}\label{sec:fp-learning-paths}

% Done
\begin{tabular}{lp{0.7\linewidth}}
    \textbf{Priorita:} & Must have \\
\end{tabular}

% Done
Aplikace by kromě kurzů měla nabízet tzv. cesty učení, neboli kariérní cesty. Tyto cesty by měly sdružovat jednotlivé kurzy do skupin podle daného tématu. Podobně jako u kurzů by každá kariérní cesta měla mít vlastní stránku s popisem a přehledem kurzů, které jsou v kariérní cestě obsaženy. Uživatel by měl mít možnost si kariérní cestu zapsat. Dále by měl mít možnost se přesunout na stránku pro správu zapsaných kurzů, kde se bude moci z kariérní cesty odhlásit. Pokud je uživatel do kariérní cesty zapsaný, měl by mít možnost na stránce cesty jednoduše pokračovat v učení dané cesty. Uživatel by měl mít také možnost zobrazit svůj certifikát o dokončení kariérní cesty, pokud jej získal.

% Done
Tento požadavek vychází z analýzy konkurence, kde většina analyzovaných aplikací nabízela určitou formu kariérních cest, díky kterým se uživatel mohl snadno naučit témata týkající se konkrétní profese.

% Done
\subsubsection{FP-9 Studijní plán}\label{sec:fp-study-plan}

% Done
\begin{tabular}{lp{0.7\linewidth}}
    \textbf{Priorita:} & Must have \\
\end{tabular}

% Done
Aplikace by měla obsahovat stránku studijního plánu uživatele, na které bude zobrazen seznam všech zapsaných kurzů uživatele. Na této stránce by měl mít uživatel možnost si ze zapsaných kurzů zvolit, které kurzy se chce učit, které si chce pouze opakovat a které se nechce momentálně vůbec učit. Podle studijního plánu pak bude aplikace doporučovat uživateli lekce při učení a opakování.

% Done
Funkce studijního plánu vychází z analýzy konkurenční aplikace Codecademy, která umožňuje uživateli si vygenerovat studijní plán. Funkce dále vychází z problému, který nastává v konkurenčních aplikacích, pokud se uživatel chce učit více kurzů zároveň. Uživatel při učení více kurzů bývá nucen často přepínat mezi jednotlivými kurzy a sám si plánovat, kterou látku se bude dále učit. Studijní plán by měl tomuto problému předejít a umožnit uživateli se snadno učit látku napříč mnoha kurzy.

% Done
\subsubsection{FP-10 Učení se a opakování podle studijního plánu}\label{sec:fp-study-plan-learning}

% Done
\begin{tabular}{lp{0.7\linewidth}}
    \textbf{Priorita:} & Must have \\
\end{tabular}

% Done
Aplikace by měla uživateli na hlavní stránce zobrazit nadcházející lekci, kterou by se uživatel měl dle jeho studijního plánu učit. Lekce by měla být vybrána ze zapsaných kurzů uživatele, které jsou ve studijním plánu v režimu učení. Uživatel by měl mít možnost si tuto lekci spustit.

% Done
Dále by měl uživatel mít možnost si na hlavní stránce aplikace spustit opakování, v rámci kterého mu aplikace vybere látku k zopakování podle toho, jak dobře si uživatel danou látku pamatuje. Látka k opakování by měla být vybrána z kurzů, které jsou ve studijním plánu buď v režimu učení, nebo v režimu opakování.

% Done
Dále by měl mít uživatel možnost si snadno spustit učení konkrétního kurzu, pokud se zrovna nechce učit doporučenou lekci podle studijního plánu.

% Done
Tato hlavní stránka by měla být nastavena jako výchozí stránka aplikace, na kterou bude uživatel přesměrován po spuštění aplikace.

% Done
Cílem této funkce je poskytnout uživateli co nejjednodušší způsob spuštění učení a využít tak princip návrhu uživatelského rozhraní \textit{microbreaks} \cite{Tidwell2020DesigningInterfaces}. Hlavní myšlenkou tohoto principu je přizpůsobit aplikaci tak, aby ji mohl uživatel použít kdykoliv, kdy má v průběhu dne jen pár minut volného času a získat tak hodnotu, kterou aplikace nabízí. Když uživatel například čeká chvíli na zastávce, může aplikaci na pár minut zapnout a využít tak efektivně svůj čas. Klíčovým prvkem tohoto principu je umožnit uživateli se ihned po spuštění aplikace přesunout k dané hodnotě, kterou aplikace nabízí. V případě této aplikace je tak snaha umožnit uživateli se ihned po spuštění aplikace přesunout k samotnému učení. Tento princip je využíván například aplikacemi sociálních sítí, které jsou pečlivě navrženy tak, aby je mohl uživatel kdykoliv v průběhu dne na pár minut zapnout a použít.

% Done
\subsubsection{FP-11 Opakování po dokončení lekce}\label{sec:fp-post-lesson-review}

% Done
\begin{tabular}{lp{0.7\linewidth}}
    \textbf{Priorita:} & Must have \\
\end{tabular}

% Done
Zásadní funkcí této aplikace je poskytnout uživateli možnost si jednoduše opakovat již naučenou látku tak, aby ji nezapomněl. Po dokončení každé nové lekce by tak měla být uživateli nabídnuta možnost opakování již naučené látky. Uživatel bude moci opakování odmítnout a přesunout se tak na stránku ze které spustil učení nové lekce. Dále bude mít možnost spustit opakování.

% Done
Po spuštění opakování bude uživateli vybrána látka k zopakování. Tato látka by měla být vybrána z kurzů, které jsou ve studijním plánu buď v režimu učení, nebo v režimu opakování.

% Done
\subsubsection{FP-12 Nová struktura učení a opakování}\label{sec:fp-learning-structure}

% Done
\begin{tabular}{lp{0.7\linewidth}}
    \textbf{Priorita:} & Must have \\
\end{tabular}

% Write - do use casů vytvořit diagram aktivit pro průchod učení

% Done
Učení by mělo být strukturováno tak, aby uživateli byla jednak po malých krocích vysvětlena probíraná látka a jednak aby si uživatel mohl danou látku procvičit. Učení se tak bude skládat z různých bloků, přičemž některé budou zaměřené na vysvětlení látky a některé na procvičení dané látky. Při procvičování látky by měla aplikace sbírat data o tom, jak uživatel danému tématu rozumí. Tato data budou později využita pro doporučování látky k opakování.

% Done
Při učení by měla aplikace uživateli vysvětlovat látku formou připomínající video s audio nahrávkou. V rámci cvičení by pak měla aplikace zobrazovat cvičení, která opět mohou být doplněna o audio nahrávku. Cvičení i vysvětlení látky budou doplněny o různé vizuální komponenty, které pomohou uživateli lépe porozumět dané látce. Učení tak bude kombinovat přístup konkurenční aplikace Mimo, která využívá různé komponenty pro vysvětlení látky, a přístup konkurenční aplikace Scrimba, v níž je látka vyučována pomocí interaktivního prostředí s audio nahrávkou vysvětlující danou látku.

% Done
Uživatel by měl mít při učení možnost se vracet k již dokončeným blokům, tedy k dokončeným cvičením a k vysvětlením látky.

% Done
Tato struktura učení se zásadně liší od struktury učení v původní aplikaci, kde probíraná látka nebyla uživateli při učení vysvětlena, učení nebylo doplněno o žádné animace nebo audio nahrávky a uživatel neměl možnost se vracet k již dokončeným blokům.

% Done
Při řešení tohoto požadavku je třeba také brát v potaz výsledky uživatelského testování stávajícího řešení, které bylo provedené v rámci mé bakalářské práce \cite{Jenicek2024BachelorThesis}. Výsledky tohoto testování ukázaly, že aplikace nedává uživateli dostatečně jasnou zpětnou vazbu o tom, zda-li odpověděl při plnění cvičení správně nebo špatně. Proto by měla aplikace uživateli také zřetelněji poskytovat zpětnou vazbu při cvičeních.

% Done
\subsubsection{FP-13 Bloky cvičení}\label{sec:fp-exercise-blocks}

% Done
\begin{tabular}{lp{0.7\linewidth}}
    \textbf{Priorita:} & Must have \\
\end{tabular}

% Done
V aplikaci by v rámci učení a opakování měly být uživateli zobrazovány bloky cvičení, díky kterým si bude moci procvičit a zopakovat danou látku. Aplikace by měla obsahovat různé komponenty, ze kterých se budou jednotlivé bloky cvičení skládat. Tyto komponenty lze rozdělit do dvou kategorií.

% Done
První kategorií jsou pasivní komponenty, které mají za cíl pouze zobrazovat informace ve formě textu, obrázků, tabulek, diagramů apod. Tyto komponenty budou sloužit zejména pro specifikaci zadání daného cvičení. Aplikace by měla obsahovat následující pasivní komponenty.

% Done
\begin{itemize}
    \item Nadpis
    \item Odstavec textu
    \item Jednoduchá tabulka
    \item Tabulka připomínající tabulku v databázi
    \item Seznamy číslované a nečíslované
    \item Obrázky ve formátech PNG, JPG, SVG a GIF
    \item Diagramy definované pomocí jazyka Mermaid\footnote{Dostupné na: \href{https://mermaid.js.org/}{https://mermaid.js.org/}}.
    \item Ukázka kódu
    \item Ukázka výstupu kódu v terminálu
    \item Ukázka výstupu kódu v podobě webové stránky
    \item Ukázka výstupu kódu v podobě tabulky
\end{itemize}

% Done
Druhou kategorií jsou aktivní komponenty, se kterými bude uživatel moci interagovat a splnit tak dané cvičení. Cílem těchto komponentů je získat data o tom, jak uživatel danému tématu rozumí. Aplikace by měla obsahovat následující aktivní komponenty.

% Done
\begin{itemize}
    \item Otázka s výběrem správné odpovědi z nabídky odpovědí
    \item Vyplnění chybějícího textu v kódu s vynechanými políčky
    \item Vyplnění chybějícího textu v kódu s vynechanými políčky pomocí nabídky slov
    \item Psaní nebo přepisování kódu ručně
\end{itemize}

% Done
U komponent, v nichž uživatel může psát celý kód, by měla aplikace umožnit uživateli jednoduše zformátovat daný kód. Tato funkce vychází z analýzy konkurence Codecademy, která tuto funkci nabízí a může přispět k lepší uživatelské zkušenosti.

% Done
Aplikace by měla umožňovat k jednotlivým komponentům připojit audio nahrávku, která bude daný komponent komentovat a doprovázet tak uživatele při vyplňování cvičení. Při zobrazení bloku cvičení by se měly postupně spustit všechny audio nahrávky, které takto provedou uživatele celým cvičením. Uživatel by měl mít možnost si audio kdykoliv vypnout a používat aplikaci i bez něj.

% Done
\subsubsection{FP-14 Editor kódu}\label{sec:fp-code-editor}

% Done
\begin{tabular}{lp{0.7\linewidth}}
    \textbf{Priorita:} & Should have \\
\end{tabular}

% Done
V rámci složitějších cvičení by měla aplikace obsahovat jednoduchý editor kódu, který bude umožňovat uživateli jednoduše psát a upravovat kód ve více záložkách připomínající různé soubory nebo okna.

% Done
V jednotlivých záložkách by mělo být možné uživateli zobrazovat následující komponenty.

% Done
\begin{itemize}
    \item Vyplnění chybějícího textu v kódu s vynechanými políčky
    \item Vyplnění chybějícího textu v kódu s vynechanými políčky pomocí nabídky slov
    \item Napsání nebo přepsání kódu ručně
    \item Ukázka výstupu kódu v terminálu
    \item Ukázka výstupu kódu v podobě webové stránky
    \item Ukázka výstupu kódu v podobě tabulky
\end{itemize}

% Done
\subsubsection{FP-15 Bloky vysvětlení látky}\label{sec:fp-explanation-blocks}

% Done
\begin{tabular}{lp{0.7\linewidth}}
    \textbf{Priorita:} & Must have \\
\end{tabular}

% Done
V rámci učení by měly být zobrazovány bloky vysvětlení látky, které mají za cíl uživateli danou látku krátce a srozumitelně vysvětlit před tím, než začne plnit jednotlivá cvičení. Celé vysvětlení látky bude probíhat formou audio výkladu, doplněného o různé animace, vizualizace a další komponenty. Tyto komponenty budou sloužit jednak jako nástroj pro lepší porozumění látky a jednak jako způsob zachycení a udržení pozornosti uživatele během výkladu. Komponenty by tak, na rozdíl od komponent bloků cvičení, měly obsahovat i různé animace a přechody tak, aby byl výklad vizuálně poutavý.

% Done
V rámci bloků vysvětlení látky by mělo být možné uživateli zobrazovat následující komponenty.

% Done 
\begin{itemize}
    \item Text s možností zvýraznění určitých slov a s možností vypsání určitých slov ve stylu kódu. % ✅
    \item Jednoduchá tabulka % ✅
    \item Tabulka připomínající tabulku v databázi % ✅
    \item Seznamy číslované a nečíslované % ✅
    \item Obrázky ve formátech PNG, JPG, SVG a GIF % ✅
    \item Diagramy definované pomocí jazyka Mermaid\footnote{Dostupné na: \href{https://mermaid.js.org/}{https://mermaid.js.org/}} s animovanými přechody mezi různými kroky úprav diagramu. Aplikace by měla podporovat diagram tříd, ER diagram, stavový diagram, diagram architektury, diagram git větví a vývojový diagram. % ✅
    \item Ukázka kódu s animovanými přechody mezi různými kroky upravování kódu. % ✅
    \item Ukázka terminálu s animovanými přechody mezi různými kroky zadávání příkazů do terminálu a vypisování jejich výstupů. % ✅
    \item Ukázka webové stránky % ✅
    \item Pozadí ve formě jednobarevného pozadí % ✅
    \item Pozadí ve formě obrázku % ✅
    \item Pozadí ve formě videa % ✅
\end{itemize}

% Done
Výklad látky by mělo být možné ovládat následujícími způsoby. Uživatel by měl mít možnost výklad zastavit, přeskočit na libovolnou část výkladu a nebo výklad zrychlit. Po skončení výkladu by měl mít možnost výklad spustit znovu.

% Done
Jelikož se může stát, že uživatel nemá vždy možnost poslouchat audio výklad, měl by mít možnost si zvuk vypnout. V takovém případě by mu měly být zobrazeny titulky, které budou vysvětlovat látku místo audio výkladu.

% Done
Zásadní vlastností celého bloku vysvětlení látky je, že by měl být responzivní a mělo by tak být možné výklad jednoduše sledovat jak na malých, tak na velkých displejích. Dále by měl výklad brát v potaz světlý resp. tmavý režim aplikace a zobrazovat všechny komponenty v příslušném režimu. 

% Done
\subsubsection{FP-16 Vysvětlení odpovědí ve cvičeních}\label{sec:fp-answer-explanation}

% Done
\begin{tabular}{lp{0.7\linewidth}}
    \textbf{Priorita:} & Should have \\
\end{tabular}

% Done
Po tom, co uživatel odpoví v rámci cvičení, by mu měla aplikace poskytnout možnost si zobrazit vysvětlení dané odpovědi. Díky tomu bude moci získat podrobnější vysvětlení, proč byla jeho odpověď ve cvičení špatná nebo správná.

% Done
\subsubsection{FP-17 Nahlašování problému při učení}\label{sec:fp-report-learning-issue}

% Done
\begin{tabular}{lp{0.7\linewidth}}
    \textbf{Priorita:} & Should have \\
\end{tabular}

% Done
Uživatel by měl mít při učení možnost u každého cvičení nebo vysvětlení látky nahlásit problém, který při dané aktivitě nastal. Při nahlašování by měl mít možnost zvolit o který typ problému z následujícího výběru se jedná.

% Done
\begin{itemize}
    \item Uživatel považuje svou odpověď ve cvičení za správnou, i přes to, že ji aplikace označila jako špatnou.
    \item Uživatel považuje svou odpověď ve cvičení za špatnou, i přes to, že ji aplikace označila jako správnou.
    \item Zadání cvičení je pro uživatele matoucí nebo nejasné.
    \item Vysvětlení látky je pro uživatele matoucí nebo nejasné.
    \item Audio nahrávka chybí nebo není v pořádku.
    \item Nastala jiná chyba. V tomto případě by měl uživatel mít možnost sám uvést, o jakou chybu se jedná.
\end{itemize}

% Done
Uživatelé s rolí korektora by měli kromě předchozích možností nahlášení problému, mít také k dispozici následující možnosti.

% Done
\begin{itemize}
    \item Vysvětlení látky je třeba upravit.
    \item Strukturu cvičení je třeba upravit.
    \item Struktura lekce by měla být upravena. Tato možnost se využije například pokud v lekci není dostatek cvičení nebo vysvětlení látky.
\end{itemize}

% Done
Korektoři by také měli mít možnost u každé odpovědi detailněji popsat vlastními slovy daný problém. 

% Done
Nahlášení problému od uživatelů s rolí korektora by mělo být označeno tak, aby bylo možné jej rozlišit od nahlášení problému běžných uživatelů. Role korektora bude nastavována uživatelům mimo frontend této aplikace. Tuto roli využijí zejména lidé zodpovědní za tvorbu a korekci jednotlivých kurzů.

% Done
\subsubsection{FP-18 Statistiky o učení}\label{sec:fp-learning-stats}

% Done
\begin{tabular}{lp{0.7\linewidth}}
    \textbf{Priorita:} & Should have \\
\end{tabular}

% Done
Aplikace by měla při učení a opakování uživatele nově sbírat data o tom, jak dlouho uživatel plnil jednotlivá cvičení, jak dlouho procházel vysvětlení látky a jakou úspěšnost měl při plnění jednotlivých cvičení. Tato data by měla být odeslána do backendové části aplikace, kde budou následně zpracována a použita pro úpravu odhadů časové náročnosti jednotlivých cvičení. Touto funkcí backendové části aplikace se zabývá kolega Bc. Jakub Vondráček v jeho diplomové práci.

% Done
Uživateli by dále měla na konci učení nebo opakování aplikace zobrazit přehled se statistikami. Na této stránce by měly být zobrazeny informace o tom, kolik času se uživatel učil, jakou měl úspěšnost při plnění cvičení, kolik nových konceptů se naučil nebo kolik konceptů si zopakoval a kolik bodů XP získal.

% Done
\subsubsection{FP-17 Stav energie}\label{sec:fp-energy-status}

% Done
\begin{tabular}{lp{0.7\linewidth}}
    \textbf{Priorita:} & Should have \\
\end{tabular}

% Done
Uživatel by měl mít možnost si v aplikaci zobrazit svůj stav energie, který bude reprezentovat, kolik lekcí nebo opakování látky může uživatel ještě za daný den spustit.

% Done
Každým spuštěním lekce nebo opakování látky se stav energie sníží. Pokud uživatel všechny body energie spotřebuje, nebude mu umožněno se dále učit. Body energie se budou každý den obnovovat.

% Done
Uživatelé s premium účtem budou mít k dispozici neomezené množství energie, tedy učení nebude body energie snižovat.

% Done
Tento požadavek vychází z analýzy konkurence, kde většina aplikací obsahuje určitou podobu omezení učení pro uživatele, kteří nemají premium účet. Tato funkce slouží primárně k motivaci uživatele, aby vybudoval pravidelný návyk učení a také k tomu, aby byl uživatel motivován si zakoupit premium účet.

% Write - požadavek na nějakou vizualizaci progressu co se týče znalostí

% Review
\subsubsection{FP-18 Registrace a přihlášení pomocí Google a Apple}

% Review
\begin{tabular}{lp{0.7\linewidth}}
    \textbf{Priorita:} & Could have \\
\end{tabular}

% Review
Uživatel by měl mít možnost se registrovat a přihlásit do aplikace pomocí účtu Google nebo Apple. 

% Review
Tento požadavek vychází z analýzy konkurence, kde většina konkurenčních aplikací tuto možnost nabízí.

% Done
\subsubsection{FP-19 Stránka nastavení}\label{sec:fp-settings-page}

% Done
\begin{tabular}{lp{0.7\linewidth}}
    \textbf{Priorita:} & Must have \\
\end{tabular}

% Done
Původní aplikace obsahovala velice jednoduchou stránku nastavení, která umožňovala uživateli nastavit pouze několik základních funkcí. Jelikož bude do aplikace přidána řada nových funkcí, bude nutné tuto stránku zásadně změnit a rozšířit tak, aby umožňovala uživateli přehledně nastavovat potřebné funkce. Stránka nastavení by tak měla poskytovat uživateli následující možnosti.

% Done
\begin{itemize}
    \item Nastavení účtu uživatele, tedy uživatelského jména, jména, emailu, hesla a avatara.
    \item Nastavení režimu vzhledu aplikace dle požadavku \nameref{sec:fp-dark-mode}.
    \item Nastavení notifikací, tedy které notifikace uživatel chce dostávat a které ne.
    \item Odhlášení se ze zapsaného kurzu.
    \item Zobrazení základní nápovědy v aplikaci.
    \item Nahlášení problému nebo chyby v aplikaci.
    \item Zobrazení podmínek a dalších právních informací o používání aplikace.
\end{itemize}

% Done
\subsubsection{FP-20 Tmavý režim}\label{sec:fp-dark-mode}

% Done
\begin{tabular}{lp{0.7\linewidth}}
    \textbf{Priorita:} & Should have \\
\end{tabular}

% Done
Rozhraní aplikace by mělo být dostupné ve dvou režimech, ve světlém a tmavém. Uživatel by měl mít možnost v nastavení zvolit, zda chce používat světlý režim, tmavý režim, nebo režim dle systémového nastavení. Rozhraní aplikace by se pak mělo zobrazovat uživateli v tomto režimu.

% Review
\subsubsection{FP-21 Push Notifikace}

% Review
\begin{tabular}{lp{0.7\linewidth}}
    \textbf{Priorita:} & Should have \\
\end{tabular}

% Review
Aplikace by měla být schopna uživateli posílat push notifikace. Tyto notifikace by měly uživatele upozorňovat v následujících situacích:

% Review
\begin{itemize}
    \item Pokud se uživatel ještě neučil v daný den.
    \item Pokud se uživatel delší dobu v aplikaci nic neučil.
    \item Pokud uživateli končí předplatné aplikace a ještě ho neobnovil.
    \item Pokud uživatele začal sledovat jiný uživatel.
    \item Pokud bylo použito zmrazení streaku.
    \item Pokud byla zveřejněna aktivita jiného uživatele, kterého daný uživatel sleduje.
    \item Pokud se blíží konec týdne a uživateli hrozí propadnutí do nižší ligy.
\end{itemize}

% Write- zde nezapomenout naimplementovat tohle: Pokud uživatel nemá zapnuté notifikace z požadavku \nameref{sec:fp-push-notifications}. Už to mám zahrnuté v požadavku pro in-app ozámení, ale neimplementoval jsem zatím tu mechaniku toho kdy se do Notices přidá toto oznámení. 

% Review
V rámci této práce se budu zabývat pouze zprovozněním zobrazování push notifikací uživateli ve frontendové části aplikace. Podrobným fungováním odesílání push notifikací se bude zabývat kolega Bc. Jakub Vondráček jeho diplomové práci, která se zabývá backendovou částí této aplikace.

% Review
\subsubsection{FP-22 Připomínání učení}

% Review
\begin{tabular}{lp{0.7\linewidth}}
    \textbf{Priorita:} & Could have \\
\end{tabular}

% Write - nakonec jsem se rozhodl, že bude povolené pouze chytré připomínání a uživatel nebude mít možnost si nastavit konkrétní čas.

% Review - tohle se ještě rozhodnout, jestli tam chci mít nebo ne
Uživatel by měl mít možnost si v aplikaci nastavit připomenutí, že se má v určitý čas učit. Aplikace by pak měla uživateli posílat push notifikace, které ho na učení upozorní, pokud se ještě v daný den neučil. Tato funkce je inspirována konkurenčními aplikacemi, které umožňují uživateli si nastavit každodenní připomenutí učení.

% Review
Dále by aplikace měla uživateli umožnit zapnout "chytré připomínání". Pokud uživatel tuto funkci zapne, nebude mu aplikace připomínat učení podle nastaveného času, ale na základě předchozího používání aplikace. Tato možnost by měla být v aplikaci nastavena jako výchozí.

% Done
\subsubsection{FP-23 Systém zobrazování reklam}

% Done
\begin{tabular}{lp{0.7\linewidth}}
    \textbf{Priorita:} & Should have \\
\end{tabular}

% Done
V původní aplikaci existoval velice jednoduchý systém pro zobrazování reklam. Po úspěšném dokončení učení se uživateli zobrazila jednoduchá reklama formou banneru.  

% Done
Nově by aplikace měla umožňovat zobrazování dalších typů reklam, v závislosti na zařízení uživatele. Na mobilních zařízeních by aplikace měla primárně zobrazovat reklamy přes celou obrazovku. Na desktopových zařízeních by pak aplikace měla zobrazovat reklamy v podobě bannerů. Reklamy by měla aplikace zobrazovat v následujících případech.

% Done
\begin{itemize}
    \item Po skončení učení, ať už po úspěšném dokončení nebo po předčasném ukončení učení uživatelem.
    \item Pokud uživatel dostal odměnu s virtuální měnou a zvolil možnost získání dvojnásobku odměny při zhlédnutí reklamy.
\end{itemize}

% Done
Tento požadavek vychází z analýzy konkurenčních řešení, kde většina aplikací zobrazuje reklamy v průběhu používání aplikace.

% Review
\subsubsection{FP-24 Nový systém předplatného}

% Review
\begin{tabular}{lp{0.7\linewidth}}
    \textbf{Priorita:} & Must have \\
\end{tabular}

% Review
Ve stávající aplikaci má uživatel možnost si zaplatit předplatné a tím získat přístup k pokročilejším funkcím. Jelikož v rámci této práce bude do aplikace přidána řada nových funkcí, je třeba tomu přizpůsobit i systém předplatného. Předplatné by nyní mělo být rozděleno do následujících kategorií.

% Review - tady asi dopsat, že premium uživatelé budou mít přístup k smart reviews a odstranit chatbota
\begin{itemize}
    \item \textbf{Uživatelé bez předplatného} budou mít přístup ke všem základním kurzům a kariérním cestám v aplikaci. Budou ovšem omezeni systémem energie, který jim umožní projít pouze omezený počet lekcí a cvičení za den. V aplikaci budou těmto uživatelům zobrazovány reklamy a budou mít omezený přístup k chatbotovi. Nebudou moci získat certifikáty o dokončení kurzů.
    \item \textbf{Uživatelé s předplatným} budou mít neomezený přístup ke všem kurzům a kariérním cestám aplikace. Nebudou jim zobrazovány reklamy, budou moci využívat chatbota s nižšími omezeními a budou moci získávat certifikáty o dokončení kurzů.
\end{itemize}

% Review
V rámci této funkce bude dále třeba zajistit lepší uživatelskou zkušenost a informovanost uživatelů, které funkce jsou dostupné v premium verzi aplikace. Pokud uživatel například chce využít funkci, která je dostupná pouze uživatelům s předplatným, měla by mu být zobrazena stránka, která ho bude informovat, že je daná funkce dostupná pouze v premium verzi aplikace.

% Review
\subsubsection{FP-25 Zkušební verze předplatného}

% Review
\begin{tabular}{lp{0.7\linewidth}}
    \textbf{Priorita:} & Should have \\
\end{tabular}

% Review
Aplikace by měla umožnit uživateli spustit zkušební verzi předplatného na určitou dobu zdarma. Tuto možnost by měli mít pouze uživatelé, kteří dříve zkušební verzi nikdy nevyužili.

% Done
\subsubsection{FP-26 Přizpůsobení avatara uživatele}

% Done
\label{sec:fp-avatar-customization}

% Done
\begin{tabular}{lp{0.7\linewidth}}
    \textbf{Priorita:} & Should have \\
\end{tabular}

% Done
Uživatel by měl mít možnost si přizpůsobit svého avatara, tedy postavu na profilovém obrázku, pomocí položek, které zakoupil v obchodě za virtuální měnu. To se liší od původní aplikace, kde měl uživatel pouze možnost měnit obrázek na svém profilu.

% Done
Uživatel by měl mít možnost si u svého avatara přizpůsobit barvu pleti, oči, barvu očí, výraz v obličeji, vlasy, barvu vlasů, doplňky, barvu brýlí, vousy, barvu vousů, pokrývku hlavy, oblečení a barvu pozadí profilového obrázku.

% Done
Tento požadavek vychází z principů gamifikace, konkrétně z principu vlastnictví \cite{Chou2019ActionableGamification}. Tento princip říká, že uživatelé mohou být více motivování, pokud mají pocit, že v rámci aplikace něco vlastní. Příkladem toho je právě přizpůsobování avatara, díky kterému si uživatelé mohou s postavou vytvořit pouto a zvýšit tak svoji motivaci pro používání aplikace. Přizpůsobování avatara tak může mít za následek častější a pravidelnější používání aplikace \cite{Birk2018CombatingAttrition}. Tento princip například využívá i velice známá vzdělávací aplikace Duolingo\footnote{Dostupné na: \href{https://www.duolingo.com/}{https://www.duolingo.com/}}, která podobnou funkci přizpůsobení avatara uživatelům nabízí.

% Done
\subsubsection{FP-27 Rozšíření obchodu a odemykání položek}\label{sec:fp-shop-unlock-items}

% Done
\begin{tabular}{lp{0.7\linewidth}}
    \textbf{Priorita:} & Could have \\
\end{tabular}

% Done
V aplikaci již existuje stránka obchodu, v němž si uživatel může zakoupit pomocí virtuální měny profilové obrázky. Vzhledem k novým požadavkům aplikace je třeba tento obchod rozšířit tak, aby bylo možné si zakoupit i položky, které jsou popsány v požadavku \nameref{sec:fp-avatar-customization}.

% Done
Dále by v obchodu měla být zavedena nová funkce odemykání položek. Některé položky by měly být dostupné ke koupi pouze pokud uživatel dosáhl určitého ocenění v aplikaci. Pokud uživatel ještě nedosáhl daného úspěchu, měly by se položky uživateli zobrazovat jako zamčené.

% Done
U zamčených položek by měl uživatel mít možnost si zobrazit informace o tom, co musí splnit, aby danou položku odemknul.

% Done
\subsubsection{FP-28 Uživatelský profil}
\label{sec:fp-user-profile}

% Done
\begin{tabular}{lp{0.7\linewidth}}
    \textbf{Priorita:} & Should have \\
\end{tabular}

% Done
V aplikaci se již nachází stránka s uživatelským profilem, která obsahuje základní informace o daném uživateli. Po zavedení nových funkcí do aplikace je třeba tuto stránku aktualizovat  a upravit její vzhled tak, aby kromě již zobrazovaných informací o uživateli přehledně zobrazovala i následující informace.

% Done
\begin{itemize}
    \item Avatar uživatele, kterého si bude uživatel moci přizpůsobit dle požadavku \nameref{sec:fp-avatar-customization}.
    \item Liga, ve které se uživatel nachází, podle požadavku \nameref{sec:fp-leaderboard}.
    \item Krátký seznam zapsaných kurzů uživatele s možností zobrazení kompletního seznamu zapsaných kurzů. Společně s kurzy by měl být zobrazen i pokrok uživatele v rámci těchto kurzů.
    \item Krátký seznam ocenění uživatele s možností zobrazení kompletního seznamu ocenění.
\end{itemize}

% Done
\subsubsection{FP-29 Zmrazení streaku}\label{sec:fp-streak-freeze}

% Done
\begin{tabular}{lp{0.7\linewidth}}
    \textbf{Priorita:} & Should have \\
\end{tabular}

% Done
Aplikace uživateli zobrazuje tzv. streak, neboli skóre, které reprezentuje počet dní v řadě, v nichž uživatel splnil alespoň jedno učení nebo opakování.

% Done
Uživatel by nově měl mít v obchodě s virtuální měnou možnost si zakoupit do zásoby několik položek tzv. zmrazení streaku. Pokud uživatel bude mít zakoupenou jednu nebo více těchto položek, a nesplní daný den žádné učení nebo opakování, jeho streak se tím nezruší a jedna tato položka se spotřebuje. Díky tomu bude uživatel moci zachovat svůj streak, i když se některý den zapomene učit.

% Done
\subsubsection{FP-30 Rozšířená oznámení v aplikaci} 

% Done
\begin{tabular}{lp{0.7\linewidth}}
    \textbf{Priorita:} & Should have \\
\end{tabular}

% Done
V původní aplikaci existoval systém, který umožňoval uživateli zobrazovat oznámení při otevření hlavní stránky aplikace. Tento jednoduchý systém z důvodu přidání nových funkcí bude třeba rozšířit tak, aby uživateli oznámil, pokud v aplikaci nastaly následující události.

% Done
\begin{itemize}
    \item Pokud aplikace automaticky použila zmrazení streaku z požadavku \nameref{sec:fp-streak-freeze}.
    \item Pokud aplikace resetovala streak uživatele.
    \item Pokud uživatel dosáhl nové významné hranice streaku, například 100 dní.
    \item Pokud byl uživatel přesunut do vyšší ligy z požadavku \nameref{sec:fp-leaderboard}.
    \item Pokud byl uživatel přesunut do nižší ligy z požadavku \nameref{sec:fp-leaderboard}.
    \item Pokud uživatel dosáhl nového úspěchu.
    \item Pokud uživatel získal truhlu s odměnou z požadavku \nameref{sec:fp-treasure-chest}.
    \item Pokud uživatel nemá zapnuté notifikace z požadavku \nameref{sec:fp-push-notifications}.
    \item Pokud by se uživateli měla zobrazit reklama z požadavku \nameref{sec:fp-advertisement}.
\end{itemize}

% Done
\subsubsection{FP-32 Truhla s odměnou}\label{sec:fp-treasure-chest}

% Done
\begin{tabular}{lp{0.7\linewidth}}
    \textbf{Priorita:} & Should have \\
\end{tabular}

% Done
Po dokončení učení by se uživateli měla zobrazit stránka s truhlou, kterou bude moci uživatel otevřít. Po otevření truhly by měl uživatel získat náhodné množství virtuální měny, jako odměnu za jeho učení. Virtuální měnu bude moci následně využít k nákupu položek v obchodě dle požadavku \nameref{sec:fp-shop-unlock-items}.

% Done
Po otevření truhly by měla být uživateli nabídnuta možnost získat dvojnásobek této odměny, pokud zhlédne reklamu.

% Done
\subsubsection{FP-31 Žebříček uživatelů}\label{sec:fp-leaderboard}

% Done
\begin{tabular}{lp{0.7\linewidth}}
    \textbf{Priorita:} & Could have \\
\end{tabular}

% Done
V aplikaci by měly být zavedeny žebříčky uživatelů, které by měly fungovat následujícím způsobem. Každý uživatel bude přiřazen do tzv. ligy učení. Každá liga učení bude reprezentovat, jak často se daný uživatel učí.

% Done
Začátkem každého týdne by měl být každý uživatel zařazen do skupiny s omezeným počtem uživatelů, kteří se nachází ve stejné lize, tedy získali předchozí týden podobné množství bodů XP. V rámci této skupiny jsou uživatelé seřazeni podle počtu bodů XP, které získali v momentálním týdnu.

% Done
Na konci každého týdne by měl být uživatel přesunut do vyšší ligy, pokud se nacházel v žebříčku uživatelů na prvních několika pozicích. Pokud se uživatel nacházel v žebříčku na posledních několika pozicích, měl by být přesunut do nižší ligy. Pokud se uživatel nacházel v žebříčku mezi těmito pozicemi, měl by zůstat v aktuální lize. Uživatel si bude moci kdykoliv zobrazit svoji pozici v žebříčku.

% Done
Tento požadavek vychází z analýzy konkurenčních řešení, zejména z analýzy aplikací Mimo a Sololearn, v nichž jsou podobné ligy implementovány.

% Review
\subsubsection{FP-34 Sdílení pokroku mezi uživateli}

% Review
\begin{tabular}{lp{0.7\linewidth}}
    \textbf{Priorita:} & Should have \\
\end{tabular}

% Review
V aplikaci existuje základní funkce pro sdílení pokroku mezi uživateli. Tato funkce je ovšem příliš jednoduchá a s implementací nových funkcí je třeba ji přizpůsobit.

% Review
Aplikace by měla být přizpůsobena na sdílení následujících zpráv mezi uživateli.

% V rámci tohoto požadavku by taky měl být odstraněn infinite scroll. Místo toho by měla aplikace jednoduše posty za posledních 7 dní a další nezobrazovat.

% Review - zde vymyslet všechny další cases
\begin{itemize}
    \item Když uživatel dosáhne určitého streaku.
    \item Když uživatel skončí na předních příčkách v lize.
    \item Když uživatel dosáhne určité úrovně v kurzu.
    \item Když uživatel dosáhne určitého ocenění.
\end{itemize}

% Write - tohle možná je nefunkční požadavek?
\subsubsection{FP-35 Odstranění přebytečných funkcí}

% Review
\begin{tabular}{lp{0.7\linewidth}}
    \textbf{Priorita:} & Must have \\
\end{tabular}

% Review
Jelikož se aplikace bude zaměřovat na výuku softwarového inženýrství, bude třeba odstranit funkce týkající se původní domény výuky jazyků. Mezi tyto funkce patří následující.

% Write - tento seznam extrahovat do kapitoly analýza stávajícího systému a zde na to pouze odkázat.
\begin{itemize}
    \item Vyhledávání slovíček
    \item Přidávání slovíček do seznamu oblíbených slovíček
    \item Cvičení zaměřená na výuku jazyků a která nejsou využitelná pro výuku softwarového inženýrství
    \item Rozdělení typů lekcí dle aspektu výuky jazyků, například gramatické lekce, lekce slovní zásoby apod.
    \item Stránka slovíčka
    \item Výběr preferovaných témat kurzu. Tato funkce v původní aplikaci sloužila jako způsob přeskočení skupin lekcí slovíček, které se netýkaly zájmů a cílů uživatele. Uživatel tak například mohl vypnutím tématu "Příroda a životní prostřední" jednoduše ze svého učení vyřadit všechny lekce týkající se tohoto tématu. V nové aplikaci ovšem bude struktura kurzů a domény změněna natolik, že tato funkce již nebude relevantní a její ponechání by mohlo mít za následek značné navýšení komplexity celého systému.
    \item Přidávání lekcí do seznamu oblíbených lekcí. Žádná z analyzovaných aplikací tuto funkci nenabízí. Funkce v rámci původní domény výuky jazyků mohla být užitečná, nicméně nyní by zbytečně mohla komplikovat uživatelské rozhraní.
    \item Odebrání stránky lekce. Tato stránka dříve zobrazovala například seznam vyučovaných slovíček v dané lekci. Zároveň bylo možné na této stránce přidávat lekci do seznamu oblíbených lekcí. V rámci domény výuky softwarového inženýrství žádná z analyzovaných aplikací samostatnou stránku pro lekci neobsahuje. Nyní ovšem stránka žádné výhody nemá a pouze komplikuje uživatelské rozhraní.
    \item Vytváření vlastních lekcí slovíček. Původní aplikace umožňovala uživatelům seskupovat slovíčka v aplikaci do vlastních lekcí. Tato funkce již není relevantní z důvodu nové domény a struktury kurzů.
\end{itemize}

\subsection{Nefunkční požadavky}

\subsubsection{NP-1 Nový design aplikace}

% Review
\begin{tabular}{lp{0.7\linewidth}}
    \textbf{Priorita:} & Should have \\
\end{tabular}

% Review
Jelikož aplikace mění své zaměření z výuky jazyků na výuku softwarového inženýrství, je třeba změnit a přizpůsobit i samotný vzhled aplikace.

% Review
Měl by být vytvořen nový design uživatelského rozhraní, který bude aplikován jak na nové funkce, tak na funkce existující.

% Review
Základní komponenty používané při implementaci by dále měly být extrahovány do samostatné knihovny tak, aby bylo možné je využít i v dalších částech projektu, jako například na webových stránkách.

% Review
Tato knihovna komponentů by měla stavět na již existujících komponentech a měla by využívat již zavedené technologie, tedy knihovny React, TypeScript, Storybook a Material UI.

% Write - přidat sem požadavek pro error handling podle standardu RFC 9457

% Done
\subsubsection{NP-2 Zpřehlednění navigace v aplikaci}

% Done
\begin{tabular}{lp{0.7\linewidth}}
    \textbf{Priorita:} & Should have \\
\end{tabular}

% Done
Aplikace by měla poskytovat uživateli přehlednější a jednodušší způsob navigace mezi jednotlivými stránkami. Z výsledků uživatelského testování v mé bakalářské práci \cite{Jenicek2024BachelorThesis}, která se zabývala původní aplikací, vyplynulo, že mají uživatelé problém s navigací mezi jednotlivými stránkami aplikace, zejména mezi jednotlivými sekcemi v daném kurzu. Jelikož do aplikace bude také přidána řada nových funkcí, bude třeba aktualizovat navigaci tak, aby uživatel se uživatel snadno orientoval v aplikaci a aby mohl snadno využívat všechny dostupné funkce.

% Done
Změna navigace se bude týkat zejména hlavní navigační lišty, bočního navigačního panelu a stránky sloužící k přepínání mezi jednotlivými kurzy.

% Review
\subsubsection{NP-2 Aktualizace kódu aplikace}

% Review
\begin{tabular}{lp{0.7\linewidth}}
    \textbf{Priorita:} & Should have \\
\end{tabular}

% Review
V kódu aplikace by měly být provedeny následující změny se záměrem zvýšit jeho kvalitu a udržitelnost a zjednodušit budoucí vývoj.

% Review
\begin{itemize}
    \item Aktualizace používaných knihoven, nástrojů a jejich konfigurací.
    \item Aktualizace používaných formatterů a linterů.
    \item Dockerizace frontendové části aplikace.
    \item Zavedení pravidel a konvencí pro nástroje, které využívají umělou inteligenci pro generaci kódu, dokumentačních komentářů a testů, s cílem podpoření vývoje aplikace a zvýšení efektivity práce.
    \item Odstranění existujícího mocku API, místo kterého by měla aplikace využívat testovací prostředí poskytované backendovou částí aplikace.
\end{itemize}

% Review
Tento požadavek byl vytvořen na základě výstupů analýzy existujícího řešení.

% Review
\subsubsection{NP-3 Zavedení automatizovaných testů}

% Review
\begin{tabular}{lp{0.7\linewidth}}
    \textbf{Priorita:} & Should have \\
\end{tabular}

% Review
Do frontendové části aplikace by měly být zavedeny automatizované testy s cílem zajištění správného fungování aplikace a zvýšení její kvality. Automatizované testy by se měly zaměřit na klíčové části aplikace, kterými jsou proces učení, registrace a přihlášení uživatele, výběr a zápis kurzů, a kupování předplatného.

% Review
Tento požadavek vychází jednak z analýzy existujícího řešení a jednak z návrhu úprav do budoucna v mé bakalářské práci \cite{Jenicek2024BachelorThesis} zabývající se zmíněným řešením.

% Review
\subsubsection{NP-4 Automaticky generované třídy API}

% Review
\begin{tabular}{lp{0.7\linewidth}}
    \textbf{Priorita:} & Should have \\
\end{tabular}

% Review
V rámci kódu frontendu by měly být automaticky generovány třídy reprezentující backendové API na základě OpenAPI specifikace \cite{OpenAPISpecification} poskytované backendovou částí aplikace.

% Review
Požadavek vychází z analýzy existujícího řešení. Jeho cílem je zjednodušení implementace a zajištění konzistence mezi frontendovou a backendovou částí aplikace.

\subsubsection{NP-5 Animace a efekty}

% Review
\begin{tabular}{lp{0.7\linewidth}}
    \textbf{Priorita:} & Should have \\
\end{tabular}

% Review
Do aplikace by měly být přidány základní animace s cílem zlepšit uživatelskou zkušenost a udržet pozornost uživatelů. Uživatelské rozhraní by tak mělo působit vizuálně poutavějším dojmem a mělo by výrazněji reagovat na interakce s uživatelem.

% Review
Tento požadavek vychází zejména z analýzy konkurenčních řešení, ve které konkurenční aplikace jako Mimo nebo Scrimba využívají animace a efekty.

% Review
\subsubsection{NP-6 Sestavování a nasazování PWA}

% Review
\begin{tabular}{lp{0.7\linewidth}}
    \textbf{Priorita:} & Should have \\
\end{tabular}

% Review
Aplikace by měla být přizpůsobena tak, aby bylo možné ji sestavit a nasadit jako PWA, neboli progresivní webovou aplikaci. Tomu by měly být přizpůsobeny i devops procesy, v rámci kterých by měla být aplikace automaticky sestavována a nasazována jako PWA.

%\subsubsection{(NP-8 Sběr statistik o uživatelském chování)}
% Review - tenhle požadavek spíš dělat nebudu.. zabere to zbytečně moc času.. RESP můžu dělat třeba základní sběr statistik o tom, jak dlouho zaberou uživatelům jednotlivá cvičení, ale to bych dal rovnou do požadavku o změně učení.. DOBRÝ BY ALE BYLO, ŽE BYCH DO KAPITOLY TESTOVÁNÍ MOHL ZAPSAT I KVANTITATIVNÍ STUDII - nicméně na to bych stejně potřeboval aby to vyplnilo hodně uživatelů, takže bych stejně asi neměl žádné výsledky té studie..


